\chapter{净秽山与药}
\chapterkana{穢れの山}
\epigraphbox{“积土成山，风雨兴焉。”}{——《荀子·劝学》}

\kaishou{城外有一座山。\par
第一章量过它的轮廓，第五章说过城外只堆得起薄薄一层浮土。\par
两者对不上，这一章讲的就是这一处对不上。}

\section{城外的隆起}

第一章量边界时，静海七号东南另有一处隆起，没有算进那几片低地。它的边缘也圆，但不规则；高约二百米，底径约三公里，比周围的低地高出十倍有余。

当时把它当作一处撞击残丘，没有深究。理由很省事：月面到处都是坑，多一个突起不算异常。

现在要回头看一眼。撞击残丘的剖面有固定形状：外缘隆起，内部下陷，坑底之下壳层减薄。这一处不是。重力场上，它的下方没有负异常，反而有一小块正异常，位置正在隆起的中轴。

正异常意味着质量多出来。一座土石堆成的丘，正异常应当弥散；这一处集中在中轴，像是有个重的东西压在底下。

\section{山里存的东西}

一九八六年，苏联人到这里来过一趟。来由是别处的事：切尔诺贝利之后，莫斯科对一切「不明放射源」都敏感到极点，凡有报告说某处读数不对，都要派人去看。月球背面一处重力正异常，加上轨道器在那里测到过几次反常计数，就被列进了名单。\footnote{那一批月面材料都是美苏探测器带回的，基金会只负责编目：六十年代起两国轮着往月面送机器，东西一批接一批，谁也说不清哪一箱出自哪一次。这一段与本书其余各章的取证方式相同。}

落点是隆起北侧，两台自动钻机，设计孔深一百二十米与一百八十米。地面站的记录后来译成中文，收在基金会编目里，编号 \texttt{ФВ-1986-04}。

\begin{archive}[内部参考]{档案 № ФВ-1986-04}
十九时〇四分 一号钻机开钻。
二十二时十一分 孔深一百四十一米，遇致密层，进尺下降。
二十三时四十分 一号机外壳变色，由银白转暗红。
零时五十分 钻杆提出时断为三截，钻头缺失。
二时二十分 二号机停机，外壳同前。
取芯两节，密封。全程剂量率未见异常。
十二日后 回收钻头表面出现深褐色疏松层，触摸即碎。
十九日后 残件局部成粉。
附记 接触芯样人员一名，三日后入院。
\end{archive}

先说最实在的一条：钻头是在孔里断的。致密层没钻穿，任务就此终止，取上来的只有断掉的那一截。

那一截带上来时已经锈了。整体转成了暗红，一星一点的锈斑都没有，十二日后表面起了一层疏松，摸一下就掉渣；再七日，局部直接成粉。铁的锈不是这样走的，锈要水和氧，月面上两样都没有。

锈只是外观。真正让这件事在当年闹起来的是仪器读数。

那一截在分解的过程中放出正电子。正电子本身不稀奇，β⁺ 衰变的核素到处都有；稀奇的是能谱对不上。任何已知的 β⁺ 核素都有特征端点能量，这一处的谱形与它们一个都不合，反而与当时理论上预言的质子衰变末态吻合。\footnote{质子衰变是大统一理论的预言之一，末态含正电子。八十年代有若干实验专门找它，都没有找到。}

一九八六年秋天，这件事在物理学界传过一阵。有人写进预印本，说这是大统一理论的第一份实验证据；也有人反对，说谱形对得上不代表机制对得上。争论了不到两年，声就小了，此后没有人再提。

为什么小，档案里没有解释。能说的只有一句：那几年苏联国内的事一桩接一桩，这类报告压在最底下；等到档案开放，看过的人已经不多。今天提起这件事的几乎没有，书里写过，档案里也还在，就是没有人调。

\section{药}

堆在山上的是什么，材料只给了一个名字，而且全书只出现过一处。

静海七号出土的残简里有一支讲净秽山的处置，简文提到山中储存之物，用的字是「药」。\\footnote{简的释文与考释见《月之都——〈广寒宫残简〉辑佚与考释》论净秽山与月兔的一章。那一本处理的是简文本身，本章用的是总目与钻孔记录，可以互校。}这一支同时记着加工的次序：抽、压、成、囤，四步，每一步由谁做，都写了。

四步之中，第三步「成」最含糊。成了什么，简文没有说形状，只说「成而不坏」，能久存。

本书按「稳定的产物」来读。抽出来的东西若是活性，就不能直接堆；先加工成一种在常温下不再活动的东西，才存得住。这与第三章说过的一句话对得上：净秽卷有一条「服药者不入环」，服药与入环两件事不能同时成立。

药与不老的关系，第五章末尾已经露出过头：服者不死，也不老，代价是此后一直产生污秽。城里没有一个人会病，因为城里没有人服药；城外的山下埋着的，正是这种药。

\section{山为什么会出事}

前一九八〇年那一次，做的是再处理，也就是把药变回去。

药稳定，是有条件的。稳定指的是它不再活动，不是它不能活动。再经一次处理，它就还原为活性污秽。那一年月兔出了事（详见第九章），处置的办法之一，是就地处理山中的存货。

处理的规模与后果，第九章再讲。此处只记一条：山从此不再稳定。原本封住的活性污秽开始在山中积累，需要持续压制。压制的人在第九章出场，压制的办法也讲。

\section{为什么上月球}

到这里，前面几件事可以合起来看了。

月之民无秽，无秽即无变化：不生育，不衰老，不生病。这一点第五章推过。他们从地面上来，把城造在月面，理由之一就藏在这条性质里：月面没有生物，也就没有变化；要维持一群不变的人，最省事的地方是那块石头。

除秽是同一件事的另一半。人身上产生的、环境里产生的，一样一样抽出去，加工成稳定的产物，封进山体。城之所以能两万年不变，靠的是这套抽、压、成、囤。

所谓场，在本书只出现一次，此处点到为止：它作用的对象是人对事实的轻重，与事实本身无关；一件事若不再被觉得重要，剩下的交给时间。这条性质为何与月面有关、如何作用于地上，本书不展开，另有一本专论。

这一章的证据是三样：一处隆起的重力异常，两个钻孔的岩芯，一支残简。三样都不厚，合起来只够说一句：城外那座山没有一层是自然堆起来的，里面装的是这座城两万年抽出来的东西，加工过，封住了。

封住的东西两千年后成了什么样子，第九章再讲。
